%% kapa: exam=1 round=33 year=2022 period=1 form=A subject=civil_law subject_ko=민법 seq=15 number=15
%% source: Q-Net 감정평가사 기출문제 (cst00309) artlSeq=5209997 file=raw/1008/5209997/2022년도 제33회 감정평가사 1차시험 1교시(민법,경제학,부동산학).pdf
%% conversion: kordoc 4.24.51 -> tools/gen1.py (auto); review=proofread-sonnet; flags=none
\begin{question}{15}
甲으로부터 대리권을 수여받지 않은 甲의 처(妻) 乙은, 자신의 오빠 A가 丙에게 부담하는 고가의 외제자동차 할부대금채무에 대하여 甲의 대리인이라고 하면서 甲을 연대보증인으로 하는 계약을 丙과 체결하였다. 이에 관한 설명으로 옳은 것은? (다툼이 있으면 판례에 따름)
\begin{choices}
  \item 甲이 乙의 무권대리행위를 추인하기 위해서는 乙의 동의를 얻어야 한다.
  \item 甲이 자동차할부대금 보증채무액 중 절반만 보증하겠다고 한 경우, 丙의 동의가 없으면 원칙적으로 무권대리행위의 추인으로서 효력이 없다.
  \item 乙의 대리행위는 일상가사대리권을 기본대리권으로 하는 권한을 넘은 표현대리가 성립한다.
  \item 계약 당시 乙이 무권대리인임을 알지 못하였던 丙이 할부대금보증계약을 철회한 후에도 甲은 乙의 무권대리행위를 추인할 수 있다.
  \item 계약 당시 乙이 무권대리인임을 알았던 丙은 甲에게 乙의 무권대리행위의 추인 여부의 확답을 최고할 수 없다.
\end{choices}
\end{question}
\answer{2}
